\scanpage{12}
\unsure{groupe|gpe} $\Gamma_0$, et d'un \unsure{$\Gamma_0$-hom|$P_0$-hom} de $P_0\longrightarrow\unsure{\mathscr{J}|\mathcal{J}|\gamma}$.

Si $\Gamma_0=\{e\}$, on trouve la notion de $G$-subordination \unsure{rigidifiée|rigidifiee}, si $\Gamma_0=G(\unsure{\mathbb{Q}|Q})$, la notion de $G$-subordination tout court. On peut expliciter la notion par la donnée d'un $F$\inserted{compatible avec $\unsure{\omega|w}$,}, et d'une \deleted{\unsure{sous}} $\Gamma_0$-rigidification de $F$ i.e. un sous-\unsure{fibre|fibré} $P_0$ de $P$, qui soit un \deleted{\unsure{sous}} $\Gamma_{0_S}$-torseur pour l'action naturelle de $\Gamma_{0_S}$\,.....

Pour cette structure, il y a un topos modulaire, qui n'est autre que \deleted{\unsure{$(\mathscr{J},$a|$(\mathscr{J},$}} $(\mathscr{J},$ avec action de $\Gamma_0$ à droite). $\forall$ Le quotient $\mathscr{J}/\Gamma_0$ (avec sa structure d'espace analytique quand $\Gamma_0$ agit proprement comme groupe discret] \deleted{est} joue le role d'une \uline{variété} \uline{de \unsure{modul|moduli|module}} \uline{de niveau} $\Gamma_0$.
% archive sheet number: 11
